% 节点势的直观含义。三点网络 s->a (容量3,费用1)，a->t (2,1)，s->t (4,4)；b=(3,0,-3)。
% 最优流 x=(2,2,1)，费用8；势 pi=(4,3,0)。约化费用 cbar=c-pi_i+pi_j：s->a 0，s->t 0，a->t -2（在上界）。
% (b) 残量网络：s->a(1,1) a->s(2,-1) t->a(2,-1) s->t(3,4) t->s(1,-4)。
% 从各点再送一单位到 t 的最便宜残量路线：s: s->t 费4；a: a->s->t 费 -1+4=3；t: 0。恰为势。
\begin{tikzpicture}[>=Stealth,line width=.8pt,font=\fontsize{8}{10}\selectfont,
  v/.style={circle,draw=FigureStroke,fill=FigureBlueFill,line width=.8pt,minimum size=6.5mm,inner sep=0pt,text=FigureInk},
  e/.style={->,draw=FigureStroke,line width=.8pt,shorten >=.5pt},
  full/.style={e,draw=FigureRed,line width=1.2pt},
  r/.style={->,draw=FigureGrid,line width=.8pt,shorten >=.5pt},
  route/.style={->,draw=FigureGreen,line width=1.5pt,shorten >=.5pt},
  lab/.style={text=FigureInk,inner sep=1.2pt,font=\fontsize{7.5}{9}\selectfont},
  pot/.style={text=FigureGreen,font=\fontsize{8}{10}\selectfont},
  note/.style={text=FigureMuted,font=\FigureNoteFont\fontsize{7.5}{9}\selectfont,align=center}]
\begin{scope}
  \node[v] (s) at (0,0) {$s$}; \node[v] (a) at (1.9,1.45) {$a$}; \node[v] (t) at (3.8,0) {$t$};
  \draw[e] (s)--node[lab,above left=-1pt,align=right]{$2/3$\\$\bar c=0$} (a);
  \draw[full] (a)--node[lab,above right=-1pt,align=left,text=FigureRed]{$2/2$ 满\\$\bar c=-2$} (t);
  \draw[e] (s)--node[lab,below,align=center]{$1/4$，$\bar c=0$} (t);
  \node[pot,left=2pt] at (s.west) {$\pi_s=4$};
  \node[pot,above=1pt] at (a.north) {$\pi_a=3$};
  \node[pot,right=2pt] at (t.east) {$\pi_t=0$};
  \node at (1.9,-1.45) {(a) 最优流与节点势};
  \node[note] at (1.9,-1.95) {弧旁为流量$/$容量与约化费用};
\end{scope}
\begin{scope}[xshift=6.6cm]
  \node[v] (s) at (0,0) {$s$}; \node[v] (a) at (1.9,1.45) {$a$}; \node[v] (t) at (3.8,0) {$t$};
  \draw[r] (s) to[bend left=14] node[lab,above left=-2pt]{$1$} (a);
  \draw[route] (a) to[bend left=14] node[lab,below right=-2pt]{$-1$} (s);
  \draw[r] (t) to[bend right=14] node[lab,above right=-2pt]{$-1$} (a);
  \draw[route] (s) to[bend left=10] node[lab,above=-1pt]{$4$} (t);
  \draw[r] (t) to[bend left=10] node[lab,below]{$-4$} (s);
  \node[pot,left=2pt] at (s.west) {$4$};
  \node[pot,above=1pt] at (a.north) {$-1+4=3$};
  \node[pot,right=2pt] at (t.east) {$0$};
  \node at (1.9,-1.45) {(b) 残量网络中再送一单位到$t$};
  \node[note] at (1.9,-1.95) {弧旁为残量费用；绿色为最便宜的路线};
\end{scope}
\end{tikzpicture}
